problem:  
Let \( X \) be a compact Calabi–Yau threefold equipped with a Ricci‑flat Kähler metric \( g \) and let \( D \subset X \) be a smooth ample divisor with positive normal bundle \( N_{D/X} \).  
Choose points \(p_i \in D\) and consider the rescaled pointed manifolds  
\[
(X, \varepsilon_i^{-2} g, p_i), \quad \varepsilon_i \to 0.
\]

**Question:**  
Assume that the pointed Cheeger–Gromov limit
\[
(Y, g_\infty, p_\infty) = \lim_{i \to \infty} (X, \varepsilon_i^{-2} g, p_i)
\]
exists as a complete noncompact Ricci‑flat Kähler manifold.  

1. Is \( (Y, g_\infty) \) necessarily *asymptotically conical* (AC), with asymptotic cone isometric to the Calabi cone metric associated with the projectivized normal bundle \( \mathbb{P}(N_{D/X}) \)?  
2. If such asymptotically conical limits exist, are their asymptotic links \(S = \mathbb{P}(N_{D/X})\) (with induced Sasaki–Einstein structure) rigid, in the sense that any sufficiently small Ricci‑flat deformation of \( (Y, g_\infty) \) fixing the asymptotic cone at infinity must arise from scaling or automorphisms of \(N_{D/X}\)?  

Equivalently: Do blow‑up limits of Ricci‑flat metrics near ample divisors in compact Calabi–Yau manifolds always produce **rigid**, asymptotically conical Calabi–Yau manifolds determined solely by the projective geometry of the normal bundle to the divisor?

---

Why is it a "good" problem:  
This problem isolates a purely **analytic and geometric** question about the fine structure of Ricci‑flat Kähler metrics near divisors in a compact Calabi–Yau manifold. It avoids ill‑defined deformation‑theoretic identifications and focuses on the existence, asymptotic form, and rigidity of potential blow‑up limits.

- **Analytic frontier:** It asks whether natural geometric degenerations of Ricci‑flat metrics yield asymptotically conical Calabi–Yau metrics, generalizing the Tian–Yau and Conlon–Hein constructions to an *intrinsic blow‑up limit context*, where no prior boundary data are imposed. This is uncharted territory in geometric analysis.  
- **Geometric rigidity:** The rigidity aspect connects local analysis of Ricci‑flat metrics to global algebraic geometry—the structure of the normal bundle \(N_{D/X}\)—thus testing how local Calabi‑Yau geometry reflects complex‑algebraic data.  
- **Conceptual simplicity, profound depth:** The problem can be stated in a few lines using standard tools (Cheeger–Gromov convergence, asymptotically conical metrics), yet its resolution would require significant new analytic and geometric ideas. It represents an elegant bridge between complex differential geometry, metric limit theory, and Kähler analytic techniques, embodying the essence of a “good” research problem: simple, precise, and deeply connected to the core geometry of Calabi–Yau spaces.